%% ============================================================
%% Chapter 3 — Proposal of a Solution
%% ============================================================

\chapter{Proposal of a Solution}\label{chapter:3_Proposal}

The literature review in Chapter~\ref{chapter:2_StateOfTheArt} identified a
need for automatic Mandarin tone assessment that is evaluated on the speech of
Polish learners and judged against expert perceptual labels rather than only
against prompt transcriptions. This chapter translates that need into a concrete
research problem and defines the requirements for the experimental solution. The
proposed solution is not a complete end-user pronunciation-training application.
It is the assessment component that such an application would require. At the
system level, pronunciation assessment is decomposed into two independent
classifiers, one for segmental phonetic accuracy and one for tonal accuracy. The
implemented experiments focus on the tonal classifier, which receives a learner
utterance and predicts the expert category perceived by a qualified Mandarin
phonetics expert.

\section{Problem Definition}
\label{sec:proposal:problem}

The central problem addressed in this thesis is the automatic classification of
Mandarin lexical tones in speech produced by Polish learners. For a given audio
recording \(x\), the system must predict one class from
\(\mathcal{T}=\{T_1,T_2,T_3,T_4,T_{\mathrm{NS}}\}\), where
\(T_1\)--\(T_4\) denote the four full lexical tones of Standard Mandarin and
\(T_{\mathrm{NS}}\) denotes a non-standard learner production that the expert did
not perceive as any well-formed Mandarin tone. The target label \(y\) is the
expert annotation of the learner's realised tone, not merely the lexical tone
expected from the written prompt. The learning task can therefore be expressed as
\begin{equation}
  \hat{y}=f_{\theta}(x), \qquad \hat{y},y \in \mathcal{T},
  \label{eq:proposal:classification_problem}
\end{equation}
where \(f_{\theta}\) is a trainable model and \(\hat{y}\) is evaluated by its
agreement with \(y\).

This formulation differs from standard native-speaker tone recognition in three
ways. First, learner speech contains productions that may fall between canonical
tone categories, especially for learners whose first language does not use tone
lexically. Studies of non-tonal-L1 learners show that tonal categories are
difficult to acquire and that the T2/T3 contrast is particularly problematic
\cite{hao2012}. Polish learners have also been shown to differ from native
Mandarin speakers in pitch height and contour shape across tones
\cite{chu2024statistical,zajdler2021}. Second, the relevant assessment target is
perceptual. In a CAPT setting, the system must approximate how an instructor
would judge the production, because feedback based only on the intended prompt
tone would ignore actual learner errors. Third, the task must handle limited
annotated learner data. This constraint favours methods that can use prior
knowledge from large speech corpora, such as Wav2Vec~2.0 and HuBERT
\cite{baevski2020,hsu2021}.

The thesis evaluates this problem through a comparative experimental study. A
ToneNet-style convolutional neural network is included as a conventional
spectrogram-based reference because CNN tone classifiers have performed well on
controlled Mandarin syllable data \cite{gao2019tonenet}. The main proposed
direction, however, is to adapt self-supervised speech models to the learner
tone task. Prior work indicates that self-supervised representations encode
Mandarin suprasegmental information and that language-matched pre-training can
improve tonal representation quality \cite{yuan2023,delafuente2024}. The study
therefore compares whether such representations are more suitable for variable
learner speech than a classifier trained from spectrogram images alone.

The following research questions guide the solution:

\begin{enumerate}[label=\textbf{RQ\arabic*.}, leftmargin=1.5cm]
  \item Which of the evaluated model families, CNN or self-supervised speech
  models, provides stronger agreement with expert tone labels on Polish learner
  speech?
  \item Does intermediate adaptation on native Mandarin speech improve
  Wav2Vec~2.0 performance on learner tone classification?
  \item Do Mandarin-specific pre-trained checkpoints improve the classification
  of learner tones compared with an English-pretrained Wav2Vec~2.0 checkpoint?
  \item Which expert categories remain most difficult after fine-tuning, and what do
  these errors imply for future CAPT use?
  \item How does evaluating against expert perceptual labels differ from the
  intended-tone evaluation that dominates the published Mandarin tone literature,
  and what does this entail for interpreting the reported scores in a CAPT
  context?
  \item Under what conditions can the tonal-accuracy classifier developed in
  this thesis be combined with an independently trained phonetic-accuracy
  classifier~\cite{brzeski2025mprp} to form a two-dimensional CAPT feedback
  signal, and what additional information (e.g.\ confidence calibration,
  prompt-tone reference) would such an integration require?
\end{enumerate}

\section{Proposed Solution Overview}
\label{sec:proposal:overview}

The proposed solution is a comparative tone-assessment pipeline composed of
three elements: a labelled learner corpus, a set of model configurations, and an
expert-aligned evaluation protocol. The learner corpus provides recordings from
Polish speakers of Mandarin together with tone annotations assigned by a
qualified expert. The model configurations represent two methodological
families: a CNN baseline using low-frequency mel-spectrogram images and five
self-supervised speech systems using Wav2Vec~2.0 or HuBERT representations. The
evaluation protocol compares the models on held-out learner recordings using
accuracy, macro-F1, weighted F1, class-level metrics, and confusion matrices.

The CNN branch follows the tone-classification tradition in which an acoustic
representation is treated as an image. Each utterance is converted into a
low-frequency mel-spectrogram and passed to a convolutional classifier. This
branch establishes whether a compact spectrogram model is sufficient for the
available learner data and provides a reference point for the self-supervised
systems.

The Wav2Vec~2.0 branch tests three adaptation routes. The first route fine-tunes
an English-pretrained Wav2Vec~2.0 Large model directly on learner recordings.
The second route inserts an intermediate native Mandarin stage using AISHELL-1
\cite{bu2017aishell}; three frame-label definitions are tested to examine where
tone supervision should be placed in the syllable. The third route starts from a
Mandarin-specific Wav2Vec~2.0 checkpoint released from WenetSpeech pre-training
\cite{zhang2022wenetspeech,tencentgamemate2022}. These routes test whether the
benefit comes from self-supervised modelling in general, from native Mandarin
adaptation, or from Mandarin-matched pre-training.

The HuBERT branch evaluates a second self-supervised backbone with two heads.
One system uses frame-level cross-entropy and majority voting, matching the
general structure of the Wav2Vec~2.0 learner systems. The second uses an
utterance-level attention pooling head trained with focal loss. This comparison
tests whether pooling over the speech frames can improve robustness to class
imbalance and to uncertain frame-level supervision.

Across all self-supervised systems, the intended output is a single expert-tone
label per utterance. The system is therefore designed as an assessment model,
not as a full pronunciation tutor. It does not generate corrective phonetic
instructions, visual pitch traces, or learner-facing feedback messages. Those
functions depend on the reliability of the classifier and are therefore treated
as possible downstream applications rather than as part of the present thesis.

\section{Requirements Specification}
\label{sec:proposal:requirements}

The requirements are divided into data, modelling, evaluation, and
reproducibility requirements.

\subsection{Data requirements}
\label{sec:proposal:requirements:data}

The learner dataset must contain audio recordings from Polish learners of
Mandarin and expert tone annotations for the same utterances. The annotation set
must cover the four full lexical tones and the non-standard learner-production
category \(T_{\mathrm{NS}}\). The linguistic neutral tone is not a target category
in the learner vocabulary used here. Because the goal is to estimate performance
on unseen learners, the main self-supervised evaluation must use
speaker-disjoint training, development, and test partitions. No speaker may
appear in more than one partition. Audio items that lack an expert label or
cannot be aligned reliably must be excluded from the self-supervised frame-level
experiments, since these systems require a speech mask for training and
aggregation.

For the intermediate Wav2Vec~2.0 route, a native Mandarin corpus must be
available for frame-level adaptation. AISHELL-1 satisfies this requirement
because it is a large, open Mandarin read-speech corpus and has been widely used
as an ASR benchmark \cite{bu2017aishell}. Its role in this thesis is not to
replace learner data but to test whether native Mandarin tone supervision
improves later transfer to Polish learner speech.

\subsection{Model requirements}
\label{sec:proposal:requirements:model}

The model set must include at least one conventional non-SSL baseline and
several self-supervised alternatives. The baseline must operate without large
pre-trained speech representations, so that the comparison can distinguish the
effect of SSL pre-training from the effect of the learner dataset alone. The SSL
models must be fine-tuned under comparable learner-data conditions and must
produce one of the five expert categories for each utterance.

The proposed model set consists of six complete setups:

\begin{enumerate}[leftmargin=1.2cm]
  \item a ToneNet-style CNN trained on low-frequency mel-spectrogram images;
  \item Wav2Vec~2.0 Large fine-tuned directly on learner speech;
  \item Wav2Vec~2.0 Large adapted first on AISHELL-1 tone labels and then
  fine-tuned on learner speech;
  \item Chinese Wav2Vec~2.0 Large fine-tuned on learner speech;
  \item Chinese HuBERT-Large fine-tuned with frame-level cross-entropy;
  \item Chinese HuBERT-Large fine-tuned with attention pooling and focal loss.
\end{enumerate}

This set is intentionally limited. It is broad enough to compare the main
modelling choices motivated by the literature, but small enough for all systems
to be trained, evaluated, and analysed within one thesis.

\subsection{Evaluation requirements}
\label{sec:proposal:requirements:evaluation}

Evaluation must measure agreement with expert labels. Accuracy alone is
insufficient because the learner corpus is class-imbalanced; high accuracy may
hide poor performance on less frequent or acoustically difficult classes.
Therefore, macro-F1 and weighted F1 must be reported alongside accuracy. The
evaluation must also include per-category precision, recall, F1, and confusion
matrices, with special attention to \(T_3\) and \(T_{\mathrm{NS}}\). \(T_3\) is
expected to be difficult because of its phonetic variability and known learner
confusion with rising tones, while \(T_{\mathrm{NS}}\) is expected to be
difficult because it groups acoustically heterogeneous productions that fall
outside the four canonical tone categories.

Model selection must be based only on development data. The held-out test set
must be used for the final comparison after checkpoint selection is complete.
For the SSL systems, the same speaker-disjoint learner split must be used so
that the results are directly comparable. If a baseline uses a different split
because of its original feature-generation pipeline, that split must be reported
explicitly and the interpretation must avoid treating it as a perfectly matched
comparison.

\subsection{Reproducibility requirements}
\label{sec:proposal:requirements:reproducibility}

Each experiment must save the information needed to reconstruct the reported
tables: split sizes, selected checkpoint score, test-set aggregate metrics,
class-level metrics, and the confusion matrix. Training curves should be saved
to allow inspection of convergence and checkpoint stability. These requirements ensure that the reported results are
not based on manual transcription from console output.

\section{Assumptions and Limitations}
\label{sec:proposal:assumptions}

The first assumption is that the expert annotation is an appropriate reference
for the realised learner tone. This is consistent with the purpose of CAPT,
where the practical goal is to approximate expert judgement. At the same time,
the thesis does not estimate inter-rater reliability because only one expert
annotation source is used. The reported scores should therefore be interpreted
as agreement with this expert reference, not as agreement with a multi-rater
consensus.

The second assumption is that utterance-level tone labels can be used as a
training signal for frame-level SSL models. In the implemented systems, the
learner label is applied to frames selected by a speech mask. This is a practical
approximation: it gives the model many supervised time steps from a short
utterance, but it does not identify the exact tone-bearing interval with
phonetic precision. The AISHELL-1 intermediate experiments partly address this
issue by testing alternative frame-label definitions, but they do not remove the
approximation for learner speech.

The third assumption is that the available learner recordings are sufficiently
representative of Polish learners in the studied educational context. The corpus
contains a large number of speakers, which is valuable for a speaker-disjoint
evaluation. However, it still reflects a specific recording protocol, set of
prompts, proficiency range, and annotation procedure. Generalisation to other
learner populations, spontaneous speech, or different recording environments
requires additional validation.

Several limitations follow from these assumptions. The quantitative experiments
focus on tone-category classification and do not report a full segmental
consonant--vowel model, fluency scoring, or sentence-level prosody evaluation.
The segmental classifier is therefore treated as a parallel component of the
overall CAPT design rather than as the main experimental object of this thesis.
The study also does not implement pedagogical feedback, even though feedback is
the eventual CAPT use case. Furthermore, the CNN baseline and the SSL systems
use different feature pipelines and, in the final implementation, different
split protocols. The CNN result is therefore best interpreted as a ToneNet-style
reference rather than as a strictly controlled architectural ablation against the
SSL models. Finally, the experiments do not train a new self-supervised model
from scratch; they fine-tune existing checkpoints because the available labelled
learner data are not sufficient for large-scale pre-training.

\section{Acceptance Criteria}
\label{sec:proposal:criteria}

The proposed solution is considered successful if it satisfies the following
criteria:

\begin{enumerate}[label=\textbf{AC\arabic*.}, leftmargin=1.5cm]
  \item The learner corpus is converted into a trainable format with explicit
  tone labels, waveform paths, and speech masks for the SSL systems.
  \item The main SSL experiments use a speaker-disjoint train/development/test
  split, and all split sizes are reported.
  \item All six planned model setups are trained or otherwise evaluated under a
  documented protocol.
  \item Final test performance is reported with accuracy, macro-F1, weighted F1,
  per-category metrics, and confusion matrices.
  \item At least one SSL setup improves over the conventional CNN reference in
  class-balanced evaluation, indicating that pre-trained speech representations
  are useful for learner-tone assessment.
  \item The comparison identifies the strongest configuration and explains
  remaining class-level weaknesses rather than reporting only a single aggregate
  score.
  \item The experiment outputs are stored in machine-readable files so that the
  tables and figures in the thesis can be traced back to saved results.
\end{enumerate}

These criteria evaluate the thesis as a scientific study. They do not require a
deployable CAPT application, real-time inference, or user-interface integration.
Those are outside the scope of the present work.

\section{Use Cases}
\label{sec:proposal:usecases}

The immediate use case is research-oriented model comparison. A researcher
selects an experimental setup, trains or fine-tunes the model, and evaluates it
on held-out learner recordings. The output is a set of quantitative measures
that show how closely the model agrees with expert tone judgements. This use
case supports the main thesis objective: determining which modelling route is
most suitable for Polish learner tone classification.

A second use case is diagnostic analysis of learner-tone errors. After
inference, the predicted and expert labels can be compared in a confusion
matrix. This makes it possible to identify systematic confusions, such as
misclassification of \(T_3\) or under-recognition of the non-standard
\(T_{\mathrm{NS}}\) category. Such analysis is useful for both model development
and language pedagogy because it shows which tonal categories require more
robust automatic handling.

A third use case is preparation for a CAPT feedback module. In a future system,
a learner would record a Mandarin word or short phrase, and the classifier would
return the perceived tone category. The feedback module could then compare this
prediction with the intended prompt tone and decide whether to mark the
production as acceptable, incorrect, or uncertain. The present thesis addresses
only the classifier needed for this scenario. It does not decide how feedback
should be phrased or how uncertainty should be shown to the learner.

The final use case is instructor support. An instructor could use model outputs
as a preliminary screening tool to locate recordings that deserve closer human
inspection. This use case is appropriate only if the system is treated as an aid
and not as a replacement for expert judgement. The experimental results in
Chapter~\ref{chapter:5_Experiments} therefore report not only the best aggregate
model but also the expert categories for which automatic assessment remains weakest.

\section{Chapter Summary}
\label{sec:proposal:summary}

This chapter defined the thesis problem as expert-aligned Mandarin tone
classification for Polish learner speech. The proposed solution is a comparative
assessment pipeline that evaluates one CNN baseline and five SSL-based systems
under documented data, modelling, and evaluation requirements. The chapter also
stated the assumptions behind the use of expert labels and frame-level
supervision, identified the main limitations, and specified acceptance criteria
for judging the experimental study. Chapter~\ref{chapter:4_Design} describes
the implementation of this proposal.
